\begin{abstract}
The Discrete Cosine Transform (DCT) is a foundational operation in digital image and video compression standards such as JPEG, requiring substantial computational throughput for real-time applications. This report details the design, verification, and synthesis of a parameterized 2D DCT accelerator implemented in SystemVerilog. To systematically explore the hardware design space, four distinct microarchitectural variants were evaluated: a baseline scalar engine (D1), a deeply pipelined scalar engine (D2), an 8-way parallel engine (D3), and a fully pipelined-parallel architecture (D4). Operating on $8 \times 8$ pixel blocks with 16-bit fixed-point arithmetic, the designs were verified against a golden MATLAB reference model, consistently achieving a high image reconstruction fidelity of 51.05 dB PSNR and an SSIM of 0.9995. Synthesis demonstrates clear area-throughput tradeoffs. The baseline D1 design minimizes resource footprint, requiring only 320 ALMs and 1 DSP block. The pipelined scalar engine (D2) proves highly effective for area-constrained applications, delivering a substantial frequency and throughput boost over D1 with minimal hardware overhead. Conversely, the unpipelined parallel engine (D3) is strictly Pareto-inefficient due to severe critical path delays that degrade overall efficiency. In contrast, the high-performance D4 architecture leverages extensive pipelining and parallel multiply-accumulate (MAC) units to reach a peak operating frequency of 195 MHz and a throughput exceeding 3.04 million blocks per second, utilizing 2,600 ALMs and 8 DSP blocks. Ultimately, the results indicate that the D4 variant provides the optimal balance of throughput and area, easily satisfying the rigorous demands of real-time high-definition video processing.
\end{abstract}
